\documentclass{beamer}

% Tema e estilo
\usetheme{Darmstadt}
\usefonttheme[onlylarge]{structurebold}
\setbeamerfont*{frametitle}{size=\normalsize,series=\bfseries}

% Idioma e codificação
\usepackage[utf8]{inputenc}
\usepackage[T1]{fontenc}

% Traduções manuais (para evitar babel no beamer)
\renewcommand{\figurename}{Figura}
\renewcommand{\tablename}{Tabela}
\renewcommand{\contentsname}{Sumário}

% Pacotes adicionais
\usepackage{graphicx}
\usepackage{times}
\usepackage[algoruled,longend]{algorithm2e}
\usepackage{tikz}
\usetikzlibrary{arrows}
\usepackage{ragged2e}
\usepackage{scalefnt}
\usepackage{booktabs}
\usepackage{amsmath}

% Remover ícones de navegação
\setbeamertemplate{navigation symbols}{}

\setbeamertemplate{footline}{
  \leavevmode
  \begin{beamercolorbox}[wd=\paperwidth,ht=0.25cm,dp=0.3cm,leftskip=0.5cm,rightskip=0.5cm]{author in head/foot}
    \hbox to \paperwidth{%
      \raisebox{-0.35cm}{\includegraphics[height=0.6cm]{Fig/cefet.png}}%
      \hfill
      \raisebox{-0.1cm}{\scriptsize Slide \insertframenumber{} de \inserttotalframenumber}%
      \hspace{0.9cm}%
    }%
  \end{beamercolorbox}
}

% Título e informações
\title[Resolução do LTPMSP via SmartILS]{Abordagem Meta-Heurística para o Problema de Agendamento de Manutenção Preventiva de Longo Prazo (LTPMSP)}
\author[Apresentador]{Getúlio Santos Mendes - Orientador: Prof. Dr. Guido Pantuza}
\institute{Centro Federal de Educação Tecnológica de Minas Gerais (CEFET-MG) - Campus Divinópolis - Engenharia de Computação}
\date{\today}

\begin{document}

\begin{frame}[plain]
    \begin{center}
        \includegraphics[height=2cm]{Fig/cefetmg.jpg}
    \end{center}
    \titlepage
\end{frame}

\begin{frame}{Sumário}
    \tableofcontents
\end{frame}

\section{Introdução}

\begin{frame}{Contexto: Manutenção preventiva}
    \begin{itemize}
        \item \textbf{A Importância da Manutenção:} Falhas inesperadas em equipamentos industriais causam atrasos e grandes prejuízos financeiros.
        \item \textbf{O Problema no Chão de Fábrica:}
        \begin{itemize}
            \item Falta de agendamento otimizado gera \textbf{ociosidade} e desperdício de equipes internas.
            \item Tarefas não alocadas a tempo precisam ser \textbf{terceirizadas}, gerando custo extra elevado para as corporações.
        \end{itemize}
        \item \textbf{Objetivo Empresarial:} Minimizar custos aproveitando ao máximo a força de trabalho interna para concluir todas as manutenções necessárias no período.
    \end{itemize}
\end{frame}

\begin{frame}{O Desafio Computacional: LTPMSP}
    \begin{itemize}
        \item O \textit{Long-Term Preventive Maintenance Scheduling Problem} (LTPMSP) busca o melhor agendamento possível dessas manutenções.
        \item \textbf{Restrições e Características:}
        \begin{itemize}
            \item Janelas de tempo rígidas (data mínima de início e máxima de fim).
            \item Requisitos restritos de compatibilidade entre habilidades técnicas das equipes e necessidades das tarefas.
            \item Capacidade de horas-homem limitadas.
        \end{itemize}
        \item \textbf{Complexidade:} Classificado como NP-Difícil \cite{Brucker2007}, sua resolução por métodos exatos torna-se computacionalmente inviável no mundo real.
    \end{itemize}
\end{frame}

\section{Fundamentação Teórica}

\begin{frame}{Modelagem Matemática (Enxuta)}
    O foco do modelo matemático PLIM é punir pesadamente a terceirização ($\omega_i$) e considerar os custos fixos das equipes internas ativadas ($\gamma_k$):
    \vspace{0.3cm}
    \begin{block}{Função Objetivo}
        \[
        \text{Minimizar } f(s) = \sum_{k \in W} \gamma_k \cdot z_k + \sum_{i \in T} \omega_i \cdot \left( 1 - \sum_{k \in W_i} y_{ik} \right)
        \]
    \end{block}
    \vspace{0.1cm}
    \textbf{Variáveis Principais:}
    \begin{itemize}
        \item $z_k$: Binária, se a equipe $k \in W$ foi ativada.
        \item $y_{ik}$: Binária, se a tarefa $i \in T$ foi delegada à equipe $k$.
    \end{itemize}
    \vspace{0.1cm}
    O modelo possui mais de 15 famílias de restrições (ex: precedência, não-sobreposição, janelas de tempo limitantes).
\end{frame}

\begin{frame}{Representação da Solução}
    \begin{itemize}
        \item Para as meta-heurísticas trabalharem rápido, usamos uma \textbf{representação indireta} em formato de permutação.
        \item Uma solução $s$ não carrega horários, apenas uma \textbf{lista de prioridade} das tarefas:
    \end{itemize}
    \vspace{0.2cm}
    \begin{center}
        \LARGE
        $s = \langle \text{T}_3, \text{T}_1, \text{T}_5, \text{T}_2, \dots, \text{T}_N \rangle$
    \end{center}
    \vspace{0.2cm}
    \begin{itemize}
        \item A responsabilidade de respeitar as restrições rígidas de tempo e alocação é delegada completamente ao \textbf{Decodificador} (Algoritmo Construtivo).
    \end{itemize}
\end{frame}

\section{Metodologia}

\begin{frame}{Decodificador: Inverted JPA}
    \begin{columns}
        \column{0.55\textwidth}
        \begin{itemize}
            \item Utilizou-se a implementação existente do Inverted JPA \cite{coelho2025simheuristic} para decodificar a permutação $s$ em um cronograma viável.
            \item \textbf{Inverted:} Agrupa as tarefas \textbf{de trás para frente}, fixando o final da tarefa no final de sua janela de limite de tempo (Deadline).
            \item Se não houver equipe e máquina simultaneamente disponíveis, a tarefa falha (Custo de terceirização $\omega_i$).
        \end{itemize}
        
        \column{0.45\textwidth}
        \begin{center}
            \includegraphics[width=\textwidth]{../Imagens/inverted_jpa.png}
        \end{center}
    \end{columns}
\end{frame}

\begin{frame}{Objetivo Geral}
    \begin{itemize}
        \item Desenvolver algoritmos eficientes para o Problema de Planejamento de Manutenção Preventiva de Longo Prazo (LTPMSP) em sua forma determinística.
    \end{itemize}
\end{frame}

\begin{frame}{Objetivos Específicos}
    \begin{itemize}
        \item Revisar a literatura sobre o LTPMSP e métodos heurísticos.
        \item Implementar e testar novas buscas locais.
        \item Implementar meta-heurísticas baseadas em Busca Local Iterada (ILS).
        \item Calibrar parâmetros das meta-heurísticas utilizando a ferramenta Irace \cite{lopez2016irace}.
    \end{itemize}
\end{frame}

\begin{frame}{Meta-Heurística: SmartILS}
    \begin{itemize}
        \item O \textit{SmartILS} adapta inteligentemente a força da perturbação baseando-se no acúmulo consecutivo de falhas.
        \item Quando não ocorre melhoria após várias iterações, a intensidade de perturbação aumenta dinamicamente.
        \item Ao encontrar uma nova melhor solução, o sistema "recompensa" a busca reduzindo a perturbação para explorar os arredores do novo ótimo local.
    \end{itemize}
\end{frame}

\begin{frame}[fragile]{SmartILS: Pseudocódigo}
    \begin{center}
    \scriptsize
        \begin{algorithm}[H]
        Gerar solução inicial $s$\;
        $s^* \leftarrow s$ (Solução incumbente)\;
        \While{critérios de parada não forem atingidos}{
            Obter solução $s'$ perturbando $s^*$ com a intensidade atual\;
            Obter solução $s''$ a partir da vizinhança de $s'$ (\textit{2-opt})\;
            \If{$f(s'') < f(s^*)$}{
                $s^* \leftarrow s''$\;
                \textbf{Reiniciar} intensidade de perturbação e contador de falhas\;
            }
            \Else{
                \textbf{Incrementar} contador de falhas\;
                \If{falhas atingirem o limite}{
                    \textbf{Aumentar} intensidade da perturbação\;
                    Zerar contador de falhas\;
                }
            }
        }
        \end{algorithm}
    \end{center}
\end{frame}

\begin{frame}{Meta-Heurística: LateAcceptanceILS}
    \begin{itemize}
        \item O \textit{LateAcceptanceILS} permite explorar novas regiões do espaço de busca aceitando temporariamente soluções piores.
        \item As iterações da busca são estruturadas em \textbf{blocos}.
        \item Ao final de um bloco sem melhoras, a melhor solução encontrada durante aquele período torna-se a nova base para as perturbações, e a intensidade da perturbação é aumentada.
    \end{itemize}
\end{frame}

\begin{frame}[fragile]{LateAcceptanceILS: Pseudocódigo}
    \begin{center}
    \scriptsize
        \begin{algorithm}[H]
        Gerar solução inicial $s$\;
        $s^* \leftarrow s$; $s_{base} \leftarrow s$\;
        \While{critérios de parada não forem atingidos}{
            Obter $s'$ perturbando $s_{base}$ com a intensidade atual\;
            Obter $s''$ a partir da vizinhança de $s'$ (\textit{2-opt})\;
            \If{$f(s'') < f(s^*)$}{
                $s^* \leftarrow s''$; $s_{base} \leftarrow s''$\;
                \textbf{Reiniciar} intensidade de perturbação e contadores\;
            }
            \Else{
                Registrar melhor solução encontrada no bloco\;
                $s_{base} \leftarrow s''$ (Aceitação Tardia)\;
                \If{tamanho do bloco atingido sem melhoras}{
                    $s_{base} \leftarrow$ melhor solução do bloco\;
                    \textbf{Aumentar} intensidade da perturbação\;
                    \textbf{Reiniciar} contadores do bloco\;
                }
            }
        }
        \end{algorithm}
    \end{center}
\end{frame}

\begin{frame}{Busca Local: Movimento e Descidas 2-opt}
    \begin{itemize}
        \item O \textbf{movimento \textit{2-opt}} inverte subsegmentos contínuos da solução para gerar vizinhos, mantendo a estrutura de dependência original das tarefas de forma mais eficaz que trocas simples (\textit{swap}).
        \vspace{0.3cm}
        \item \textbf{Métodos de Descida avaliados:}
        \begin{enumerate}
            \item \textbf{Descida 2-opt Completa:} Avalia exaustivamente toda a vizinhança $O(n^2)$ iterativamente até atingir o ótimo local.
            \item \textbf{Descida 2-opt Randômica:} Reversões aleatórias sob a regra de Primeira Melhoria (\textit{First Improvement}), focando em maior velocidade.
            \item \textbf{Descida 2-opt de Segmento Randômico:} Varredura restrita apenas a blocos compreendidos entre 10\% e 30\% da permutação.
        \end{enumerate}
    \end{itemize}
\end{frame}

\begin{frame}{Calibração com Irace}
    \begin{itemize}
        \item \textbf{Espaço de Busca de Parâmetros:} A ferramenta \texttt{irace} \cite{lopez2016irace} foi usada para testar sistematicamente combinações de parâmetros entre:
        \begin{itemize}
            \item \textbf{2 Meta-Heurísticas:} \textit{SmartILS} vs. \textit{LateAcceptanceILS}.
            \item \textbf{3 Métodos de Descida 2-opt:} Completa, Randômica e Segmento Randômico.
            \item \textbf{Parâmetros Numéricos:} Limite da busca local (\texttt{itermax}), limite global (\texttt{ilsmax}) e nível de tolerância à estagnação (\texttt{nivel}).
        \end{itemize}
        \vspace{0.3cm}
        \item \textbf{Instâncias de Calibração:} Foram selecionadas \textbf{três instâncias} representativas de médio a grande porte da literatura
        \vspace{0.3cm}
        \item \textbf{Combinação vencedora:} \textit{SmartILS} Descida 2-opt Randômica, \texttt{itermax} = 5, \texttt{nivel} = 10, e \texttt{ilsmax} = 100
    \end{itemize}
\end{frame}

\section{Resultados e Conclusão}

\begin{frame}{Resultados Parciais}
    \textbf{Comparação: ILS Original vs. SmartILS}
    
    \vspace{0.2cm}
    \resizebox{\textwidth}{!}{
    \begin{tabular}{lrrrrr}
    \toprule
    \textbf{Instância} & \multicolumn{2}{c}{\textbf{Literatura (ILS Original)}} & \multicolumn{2}{c}{\textbf{Proposto (SmartILS)}} & \textbf{Gap FO (\%)} \\
    \cmidrule(lr){2-3} \cmidrule(lr){4-5}
     & \textbf{Custo Médio} & \textbf{Tempo (s)} & \textbf{Custo Médio} & \textbf{Tempo (s)} & \\
    \midrule
    LTPMSP\_01\_00300 & 1937.40 & 44.06 & 1991.10 & 0.47 & -2.77\% \\
    LTPMSP\_01\_01200 & 11081.60 & 1200.06 & 11362.80 & 38.24 & -2.54\% \\
    LTPMSP\_02\_01200 & 9510.00 & 1200.05 & 9528.60 & 27.24 & -0.20\% \\
    LTPMSP\_03\_04800 & 85037.70 & 4800.51 & 89474.40 & 2314.15 & -5.22\% \\
    LTPMSP\_06\_09600 & 30102.50 & 9602.01 & 33687.80 & 3600.18 & -11.91\% \\
    \bottomrule
    \end{tabular}
    }
\end{frame}

\begin{frame}{Resultados Parciais: Análise}
    \begin{itemize}
        \item \textbf{Speedup:} O tempo computacional apresentou uma aceleração média de \textbf{4,5 vezes}, caindo da ordem de milhares para dezenas de segundos nas maiores instâncias.
        \vspace{0.5cm}
        \item \textbf{Trade-off (Gap):} Houve um ligeiro aumento médio no custo de apenas \textbf{2,65\%} em relação à literatura original. Trata-se de um resultado satisfatório que viabiliza a execução rápida e abre espaço para melhorias futuras em cenários de incerteza.
    \end{itemize}
    \vspace{0.5cm}
    \scriptsize{* Os resultados do SmartILS refletem a média de 10 execuções independentes por instância.}
\end{frame}

\begin{frame}{Conclusões}
    \begin{itemize}
        \item \textbf{Representação e Constrição:} O uso do Inverted JPA provou ser crucial, eliminando buscas inviáveis no espaço de soluções ao embutir as restrições duras na fase de decodificação.
        \vspace{0.5cm}
        \item \textbf{Aptidão do SmartILS:} O controle inteligente da intensidade de perturbação otimizou radicalmente a métrica temporal ao evitar buscas exaustivas em vizinhanças infrutíferas.
    \end{itemize}
\end{frame}

\begin{frame}{Trabalhos Futuros (PFC II)}
    \begin{itemize}
        \item \textbf{SIM-ILS:} Desenvolver abordagem sim-heurística acoplando simulações de Monte Carlo para modelar a incerteza no tempo das tarefas.
        \vspace{0.3cm}
        \item \textbf{Vizinhança Variável (VNS e VND):} Mesclar diferentes movimentos (\textit{swap}, \textit{insertion}, \textit{2-opt}) para intensificar a busca local.
        \vspace{0.3cm}
        \item \textbf{Paralelização Heurística:} Utilizar múltiplas \textit{threads} para acelerar a avaliação simultânea dos movimentos na vizinhança.
    \end{itemize}
\end{frame}

\begin{frame}[plain]
    \begin{center}
        \LARGE \textbf{Muito Obrigado!}
        
        \vspace{1cm}

        \Large \textbf{Perguntas?}
        
        \vspace{1cm}
        \normalsize
        Contato: getuliosantosmendes@gmail.com
        
        \vspace{1.5cm}
        \includegraphics[height=2cm]{Fig/cefetmg.jpg}
    \end{center}
\end{frame}

\begin{frame}[allowframebreaks]{Referências}
    \tiny
    \bibliographystyle{apalike}
    \bibliography{../Referencias}
\end{frame}

\end{document}